\documentclass[a4paper,12pt]{article}

\usepackage{geometry}
% задание полей текста
\geometry{left=10mm,right=10mm,top=20mm,bottom=20mm}
\usepackage{listings}
\usepackage{cmap}					% поиск в PDF
\usepackage[T2A]{fontenc}			% кодировка
\usepackage[utf8]{inputenc}
\usepackage[english, russian]{babel}	% локализация и переносы
\usepackage{natbib}

\usepackage{graphicx}
\graphicspath{{pix/}}
\usepackage{pgfplots}
\usetikzlibrary{arrows}
\pgfplotsset{width=10cm,compat=newest}
\pgfkeys{/pgf/trig format=rad}


\usepackage{xcolor}
\usepackage{hyperref}
\usepackage{amsmath}
\usepackage{amssymb}
\usepackage{enumerate}
\usepackage[normalem]{ulem}


\setlength\parindent{0pt}


\sloppy                     % строго соблюдать границы текста
\linespread{1.3}            % коэффициент межстрочного интервала
\setlength{\parskip}{0.5em} % вертик. интервал между абзацами

\setcounter{secnumdepth}{0} % отключение нумерации разделов
\binoppenalty=1000          % уменьшение переносов в формулах

% объявление новых макрокоманд

\newcommand{\Def}{\textbf{Def.} }
\newcommand{\Th}{\textbf{Теорема.} }
\newcommand{\Thbd}{\textbf{Теорема (б/д).} }
\newcommand{\Theor}[1]{\textbf{Теорема ({#1}).} }
\newcommand{\Theorbd}[1]{\textbf{Теорема ({#1}) (б/д).} }
\newcommand{\Consequence}{\textbf{Следствие.} }
\newcommand{\Remind}{\textbf{Remind.} }
\newcommand{\Note}{\textbf{Note.} }
\newcommand{\Statement}{\textbf{Утверждение.} }
\newcommand{\Prop}{\textbf{Свойство:} }
\newcommand{\Props}{\textbf{Свойства:} }
\newcommand{\Proof}{\textbf{Доказательство:} }
\newcommand{\Prooff}{\textbf{Доказать:} }
\newcommand{\Solution}{\textbf{Решение:} }
\newcommand{\Alg}{\textbf{Algorithm.} }
\newcommand{\Lemma}{\textbf{Лемма.} }
\newcommand{\Example}{\textbf{Пример:} }
\newcommand{\Task}{\textbf{Задача.} }
\newcommand{\Solve}{\textbf{Решение:} }
\newcommand{\Examples}{\textbf{Примеры.} }

\allowdisplaybreaks[4]

%\renewcommand\thesection{\arabic{section}}

\newcommand{\Endproof}{$\blacksquare$ }

\newcommand{\tr}{\text{tr}}
\newcommand{\Le}{\leqslant}                % русский стиль нестрогих неравенств
\newcommand{\Ge}{\geqslant}
% облегчение математических обозначений
\newcommand{\A}{\mathcal{A}}
\newcommand{\M}{\mathcal{M}}
\newcommand{\kukarek}{\textit{(КуКаРеК) }}
\newcommand{\F}{\mathcal{F}}
\newcommand{\Gs}{\mathcal{G}}
\newcommand{\R}{\mathbb{R}}
\newcommand{\N}{\mathbb{N}}
\newcommand{\Q}{\mathbb{Q}}
\newcommand{\Norm}{\mathcal{N}}
%\newcommand{\C}{\mathbb{C}}
\newcommand{\ind}{\perp \!\!\! \perp}
\newcommand{\Z}{\mathbb{Z}}
\newcommand{\E}{{\mathbb{E}}}
\newcommand{\D} { {\mathbb{D}}}
\newcommand{\x}{\overline{x}}
\newcommand{\y}{\overline{y}}
\newcommand{\soch}[2]{\text{$C_{#1}^{#2}$}}
\newcommand{\X}{\text{$\mathcal{X}$}}
\newcommand{\VC}[2]{\text{$VC\left(#1; #2\right)$}}
\renewcommand{\P}{\text{$\mathcal{P}$}}
\newcommand{\B}{\text{$\mathcal{B}$}}
\newcommand{\brackets}[1]{\left({#1}\right)} % автоматический размер скобочек
\newcommand{\abs}[1]{\left|{#1}\right|} % автоматический размер модуля
\newcommand{\norm}[1]{\left\|{#1}\right\|} % автоматический размер нормы
\newcommand{\set}[1]{\left\{{#1}\right\}} % автоматический размер множества
\newcommand{\condset}[2]{\set{#1 \ | \ #2}} % определение множества

\newcommand{\distrfamily}{\condset{P_\theta}{\theta \in \Theta}}

\newcommand{\htheta}{\hat{\theta}}
\newcommand{\hthetax}{\htheta(X)}
\newcommand{\hthetanx}{\htheta_n(X)}
\newcommand{\stheta}{\theta^*}
\newcommand{\sthetax}{\stheta(X)}
\newcommand{\taut}{\tau(\theta)}
\newcommand{\pt}{p_\theta}
\newcommand{\Pt}{P_\theta}
\newcommand{\Et}{\E_\theta}
\newcommand{\Utx}{\text{U}_\theta(X)}
\newcommand{\Var}{\text{Var}}
\newcommand{\Vart}{\Var_\theta}

\newcommand{\samp}[2]{{#1}_1, \ldots, {#1}_{#2}}
\newcommand{\observation}{X = (\samp{X}{n})}
\newcommand{\infsamp}[1]{{#1}_1, {#1}_2, \ldots}
\newcommand{\sampf}[3]{{#1}_1(#3), \ldots, {#1}_{#2}(#3)}

\newcommand{\charac}[2]{\varphi_{#1}\brackets{#2}} % Хар.функция от распределения с.в. от аргумента
\newcommand{\proj}[2]{\text{proj}_{{#1}} {#2}}

\newcommand{\pn}{\text{п.н.}}
\newcommand{\pnt}{\text{P$_\theta$-п.н.}}
\newcommand{\pnto}{\xrightarrow{\pn}}
\newcommand{\pntto}{\xrightarrow{\pnt}}
\newcommand{\pto}{\xrightarrow{\text{P}}}
\newcommand{\ptto}{\xrightarrow{\text{P$_\theta$}}}
\newcommand{\lpto}{\xrightarrow{\text{L$_p$}}}
\newcommand{\dto}{\xrightarrow{\text{d}}}
\newcommand{\dtto}{\xrightarrow{\text{d$_\theta$}}}
\newcommand{\eps}{\varepsilon}
\newcommand{\Eps}{\mathcal{E}}
\newcommand{\probspace}{(\Omega,\: \F,\: \P)}
\newcommand{\ddfrac}[2]{\frac{\displaystyle #1}{\displaystyle #2}}

\newcommand{\todo}[1]{\section{TODO}
\begin{center}
    \LARGE #1
\end{center}}

\newcommand{\drawsome}[1]{
    \begin{figure}[h!]
            \centering
            \includegraphics[scale=0.7]{#1}
            \label{fig:first}
    \end{figure}
}

\newcommand{\drawsomemedium}[1]{
    \begin{figure}[h!]
            \centering
            \includegraphics[scale=0.45]{#1}
            \label{fig:first}
    \end{figure}
}

\newcommand{\drawsomesmall}[1]{
    \begin{figure}[h!]
            \centering
            \includegraphics[scale=0.3]{#1}
            \label{fig:first}
    \end{figure}
}


\hypersetup{
    pdfstartview=FitH,
    colorlinks=true,
    linkcolor=blue,
    filecolor=magenta,
    urlcolor=cyan,
}

\usepackage{titleps}                       % колонтитулы

\newpagestyle{main}{
  \setheadrule{.4pt}
  \sethead{\coursename}{}{\hyperlink{intro}{\;Назад к содержанию}}
  \setfootrule{.4pt}
  \setfoot{\coursedate \; МФТИ}{}{\thepage}
}

\pagestyle{main}
